\clearpage\thispagestyle{otFront}\phantomsection\label{parcours}\pdfbookmark[0]{Construire son parcours}{bm-parcours}\otPageTitre{Construire son parcours}{Avancer avec des retours réguliers}\fontsize{11}{15}\selectfont

\begin{tabularx}{\linewidth}{@{}lY@{}}\toprule
Étape & Travail conseillé\\\midrule
Consolider & Chapitres 1 à 7 : nombres, calculs, équations, proportionnalité et unités.\\[3mm]
Modéliser & Chapitres 8 à 11 : fonctions, statistiques, expériences aléatoires.\\[3mm]
Démontrer & Chapitres 12 à 16 : théorèmes, transformations, mesures et espace.\\[3mm]
Programmer & Chapitre 17, à reprendre pendant toute l’année.\\[3mm]
Synthétiser & Chapitre 18, défis et brevets blancs.\\\bottomrule
\end{tabularx}
\begin{savoirfaire}[Une séance autonome de 35 à 50 minutes]Relire une propriété et son exemple ; traiter deux exercices sans correction ; comparer la démarche ; refaire une question ratée. Terminer par une phrase : « La prochaine fois, je dois penser à… ».\end{savoirfaire}
\begin{savoirfaire}[Une révision espacée]Revenir sur une notion quelques jours après, puis quelques semaines plus tard. Mélanger un calcul, une lecture graphique et une justification géométrique évite de dépendre seulement de l’ordre du chapitre.\end{savoirfaire}
Pour un défi, prévoir 30 à 60 minutes, éventuellement en deux séances.
Pour un brevet blanc, utiliser deux copies et respecter la remise de la
première au bout de 20 minutes. Refaire ensuite les erreurs par notion,
sans attendre le prochain sujet complet.

\clearpage
